%%%%

%%%   Самарский государственный технический университет

%%%   Кафедра прикладной математики и информатики

%%%

%%%   Преамбула для статьи в журнал "Вестник Самарского государственного

%%%   технического университета. Серия: «Физико-математические науки»"

%%%   Ориентирована под MikTEx 2.4 for Win32

%%%

%%%   Хотя сам журнал собирается под GNU/Linux на дистрибутиве TeXLive



\documentclass[11pt,a4paper,twoside]{article}



\usepackage[cp1251]{inputenc}

\usepackage[T2A]{fontenc}

\usepackage[english,russian]{babel}

\usepackage{samgtubib}

\usepackage[dvips]{graphicx}

\usepackage[multiple]{footmisc}



\usepackage{samgtu}

\renewcommand{\VestnikYear}{2009} % Год издания Вестника

\renewcommand{\VestnikNumber}{2\,(19)} % Номер Вестника

\renewcommand{\VestnikMonth}{Ноябрь} % Месяц выхода Вестника

\renewcommand{\VestnikData}{01.11.09} % Дата подписания Вестника в печать

\renewcommand{\VestnikUPL}{26,64} % Усл. печ. л.

\renewcommand{\VestnikUIL}{25,73} % Уч.-изд. л.

\renewcommand{\VestnikRegNumber}{479/09} % Регистрационный номер

\renewcommand{\VestnikOrder}{994} % Номер заказа







%

%\usepackage{pst-barcode}

%\usepackage{auto-pst-pdf}

%%

%%\newcommand{\totop}[1]{\raisebox{6pt}[1pt][2pt]{#1}} 

%%\newcommand{\tottop}[1]{\raisebox{12pt}[1pt][2pt]{#1}} 

%%\newcommand{\totttop}[1]{\raisebox{18pt}[1pt][2pt]{#1}} 

%%\newcommand{\tottttop}[1]{\raisebox{24pt}[1pt][2pt]{#1}} 

%%\newcommand{\totttttop}[1]{\raisebox{30pt}[1pt][2pt]{#1}} 

%%\newcommand{\tobot}[1]{\raisebox{-6pt}[1pt][2pt]{#1}} 

%%\newcommand{\tobbot}[1]{\raisebox{-12pt}[1pt][2pt]{#1}} 

%%\newcommand{\tobbbot}[1]{\raisebox{-18pt}[1pt][2pt]{#1}}

%

%\graphicspath{{./00PIC/}}

%

%%\samgtupubl % для генерации pdf, отдаваемого в типографию СамГТУ

%\pdfcrop % обрезка pdf для размещения на math-net.ru

%\draft % На корректуре

%\filllastpage % Последняя страница не заполнена не полностью

%%\briefcomm % Статья является кратким сообщением



\vsgtu{1378}



\FirstPage{259}

\LastPage{269}

%%

%%

%\begin{document}

%



\hfill \begin{pspicture}(1in,1in)

\psbarcode{http://dx.doi.org/10.14498/vsgtu1378}{}{qrcode}

\end{pspicture}

\vspace{-1in}



% Номер УДК

\udk{539.126}%Обработка текста. Подготовка текстов : Языки программирования. Компьютерные языки

\msc{81U05, 81U10}% Коды MSC см. http://www.ams.org/msc/



\rutitle{Описание радиационных распадов $V \to  P\gamma^{*}$ в различных формах Пуанкаре-инвариантной квантовой механики}

{Описание радиационных распадов $V \to  P\gamma^{*}$ \dots}



\titlehack{Описание радиационных распадов $\boldsymbol{V \to  P\gamma^{*}}$ \\ 

в различных формах Пуанкаре-инвариантной \\

квантовой механики$^*$}



\ruauthor{А.~Ф.~Крутов, Р.~Г.~Полежаев}{К\,р\,у\,т\,о\,в~А.~Ф.,\ П\,о\,л\,е\,ж\,а\,е\,в~Р.~Г.}



\ruaffil{\rusamgu.}

\enaffil{\ensamgu.}



%\email{E-mails}{krutov@samsu.ru, polezaev@list.ru} %E-mail's авторов



\ruauthordetails{2}{\fullauthor{\rupid{20496}{Александр Федорович Крутов}}

\regalia{д.ф.-м.н., проф.; \mem{krutov@samsu.ru}}\position{профессор}

\dept{каф. общей и теоретической физики} \\ 

\fullauthor{\rupid{76307}{Роман Геннадьевич Полежаев}} \regalia{\mem{polezaev@list.ru}; автор, ведущий переписку} \position{аспирант} \dept{каф. общей и теоретической физики}

\par

\smallskip

\noindent

$^*$Настоящая статья представляет собой расширенный вариант доклада \cite{krutov:add1-ru}, сделанного авторами на Четвёртой международной конференции «Математическая физика и её приложения» (Россия, Самара, 25 августа -- 1 сентября 2014).}



\enauthordetails{2}{\fullauthor{\enpid{20496}{Alexander F. Krutov}} \regalia{Dr. Phys. \& Math. Sci.; \mem{krutov@samsu.ru}} \position{Professor} \dept{Dept. of General and Theoretical Physics}\\

\fullauthor{\enpid{76307}{Roman G. Polezhaev}} \regalia{\mem{polezaev@list.ru}; Corresponding Author} \position{Postgraduate Student} \dept{Dept. of General and Theoretical Physics}

\par

\smallskip

\noindent

$^*$This paper is an extended version of the paper \cite{krutov:add1-eng}, presented at the Mathematical Physics and Its Applications 2014 Conference.

}





\rukeywords{Пуанкаре-инвариантная квантовая механика, радиационные

распады, электромагнитный ток, переходной формфактор}



\ruabstract{Работа посвящена описанию радиационных распадов $V \to P

\gamma^{*}$ в различных основных формах  Пуанкаре-инвариантной

квантовой механики (ПИКМ). Для построения матричного элемента

электромагнитного тока перехода используется специальная

процедура, удовлетворяющая условиям лоренц-ковариантности и

сохранения. В качестве иллюстрации развитого формализма  в

модифицированном релятивистском импульсном приближении проведено

описание радиационного перехода $\rho \to \pi \gamma^{*}$.

Получено аналитическое выражение для переходного формфактора

$F_{\pi \rho}(Q^2)$, совпадающее во всех формах ПИКМ. Выполнены

численные расчеты переходного формфактора с двумя типами модельных

волновых функций.}

%

\entitle{Description of radiative decays of $V \to P \gamma^{*}$ in different forms of Poincar\'e-invariant quantum mechanics}



\engtitlehack{Description of radiative decays of $V \to P \gamma^{*}$ \\

 in different forms of Poincar\'e-invariant\\ 

 quantum mechanics$^*$}



%

\enauthor{A.~F.~Krutov, R.~G.~Polezhaev}{K\,r\,u\,t\,o\,v~A.~F., P\,o\,l\,e\,z\,h\,a\,e\,v~R.~G.}

%

%

%\enginstitute{\engsamgu.}% Место работы каждого из авторов на англ. языке

%% Если несколько, то так  \enginstitute{ Организация 1 \and Организация 2  \and Организация 3}

%



%

\enabstract{The description of the radiative decays of $V \to P \gamma^*$ in different forms Poincar\'e invariant quantum mechanics (PIQM) is considered. To construct the matrix element of the electromagnetic current we use the non-diagonal parametrization procedure. The obtained matrix element of the current satisfies the conditions of the Lorentz covariance and conservation. To illustrate this approach in a modified relativistic impulse approximation the description of the radiative transition $\rho \to \pi \gamma^{*}$ is performed. An analytic expression for the transition form factor $F_{\pi \rho}(Q^2)$, matching all forms PIQM, is obtained. Numerical calculations of the transition form factor are made with different model wave functions.}



\enkeywords{Poincar\'e invariant quantum mechanics, radiative decays, electromagnetic current, transition formfactor}



\date{17/XII/2014} % Дата поступления статьи в редакцию.

\revisiondate{12/III/2015} % В окончательном виде

\accepteddate{08/IV/2015}



%%%%%%%%%%%%%%%%%%%%%%%%%%%%%%%%%%%%%%%%%%%%%%%%%%%%%%%%%%%%%%%%%%%%%%%%%%%%%%%%%%%%





\renewcommand{\baselinestretch}{0.90}



\makerutitle



%Каждый пункт статьи должен начинаться с команды Название пункта

%после названия точку ставить не нужно. Пункты автоматически нумеруются.

%Если пункт нумеровать не нужно, то следует использовать в команде

% \Section необязательный параметр [n]:

%   \Section[n]{Имя пункта}

% Введение и заключение обычно не нумеруются.



В последние несколько лет был проведен ряд экспериментов по

изучению радиационных распадов векторных мезонов. Так

коллаборациями NA~60 [\citen{Na6009}--\citen{Na60Ur11}] и KLOE-2

\cite{Kloe11}, были измерены переходные формфакторы в реакциях

$\omega  \to  \pi \gamma^{*}$, $\phi  \to  \eta \gamma$

и определены соответствующие константы распадов. В~радиационном

переходе $\rho \to \pi\gamma^{*}$ изучались вклады пиона в

сечение распадов векторных мезонов и получены соответствующие

ширины радиационного распада \cite{Amsler08}. Указанные процессы

дают важную информацию о непертурбативной кварковой динамике на

средних и больших расстояниях.



Одним из основных теоретических методов описания перечисленных

выше процессов является Пуанкаре-инвариантная квантовая механика

(ПИКМ), которая может быть сформулирована в различных формах "---

мгновенная форма, точечная форма и динамика на световом фронте (см,

например, \cite{KrT09}). Одним из важных вопросов, возникающих при

описании составных систем, остается выбор формы ПИКМ. В работах

[\citen{MaT02}--\citen{Klink14}] при вычислении электромагнитных

формфакторов скалярных и векторных мезонов для основных форм ПИКМ

получаются разные результаты. Причина расхождения результатов,

по-видимому, связана с различием приближений, используемых при

построении матричного элемента тока и, в частности, с учетом

условий лоренц-ковариантности и сохранения \cite{KrT09}. Таким

образом, формулировка новых  способов построения операторов

электромагнитных переходов является в настоящее время актуальной

задачей.



Отметим, что мгновенная форма ПИКМ была с успехом применена

авторами для описания электромагнитной структуры таких систем с

сильным взаимодействием, как пион, $\rho$-мезон, дейтрон, включая

области больших переданных импульсов \cite{KrT09, KrTrTs09}. 

В~этих работах построение тока осуществлялось для

случаев диагональных по полному угловому моменту матричных

элементов. В настоящей работе для описания процессов радиационных

распадов векторных мезонов $V  \to  P \gamma^{*}$ проводится

обобщение методики построения тока на недиагональные по полному

угловому моменту матричные элементы и на другие формы ПИКМ.



Ядром развитого подхода является  процедура построения матричных

элементов локальных операторов (см., например, \cite{KrT09}). Эта

процедура позволяет выделить из матричного элемента оператора

любой тензорной размерности приведенные матричные элементы

(формфакторы), которые являются инвариантами при преобразованиях

из группы Пуанкаре.



В рамках рассматриваемого подхода с использованием

модифицированного импульсного приближения была показана

эквивалентность основных форм ПИКМ на примере расчета

электромагнитного формфактора пиона \cite{KrPol13-2}.



В представленной работе показана эквивалентность основных форм \linebreak

ПИКМ при расчете переходного формфактора радиационного перехода \linebreak

\mbox{$\rho  \to  \pi \gamma^{*}$}. 

Используемый при этом

математический аппарат параметризации матричных элементов

локальных операторов \cite{ShirCh63} применен для описания

недиагональных по полному угловому моменту матричных элементов

тока.



В работе произведен также расчет переходного формфактора процесса

\mbox{$\rho  \to  \pi \gamma^{*}$} с~ двумя типами модельных

волновых функций.



\smallskip



\Section{Построение матричного элемента электромагнитного тока перехода для систем

невзаимодействующих частиц с квантовыми числами $\boldsymbol \rho$- и $\boldsymbol \pi$-мезонов}

Рассмотрим параметризацию матричного элемента тока между

состояниями с квантовыми числами $\rho$- и $\pi$-мезонов, т.е.

недиагонального по полному угловому моменту:

\begin{equation}

\label{krutov:nodia}

 \langle\vec{W^{i}},\sqrt{s}|j_{\mu}^{0}(0)|\vec{W'^{i}},\sqrt{s'}, 1, 0, 1, m'\rangle\;.

 \end{equation}

Здесь

$\vec{W^{i}}=\vec{w^{i}_{1}}+{\vec{w^{i}_{2}}}$;

$\sqrt{s}$ "---

инвариантная масса системы двух свободных частиц;

$\vec{w^{i}}$ "--- трехмерный вектор, различный для разных форм ПИКМ.

В мгновенной форме $\vec{w^{1}}=\vec{p}$ "--- импульс, в точечной 

$\vec{w^{2}}=\vec{v}$ "--- 4-скорость и в динамике на световом фронте "---

$\vec{w^{3}}=\tilde{\vec{p}}$, где

 \begin{equation}

 %\label{krutov:1}

\tilde{\vec{{p}}}=(p_{\bot},p^{+}), \quad p_{\bot}=(p^{1},p^{2}),

\quad p^{+}=(p^{0}+p^{3})/\sqrt{2}\;.

\end{equation}





Для инвариантной параметризации матричного элемента \eqref{krutov:nodia} выполним преобразование Лоренца из исходной (лабораторной) системы координат в систему Брейта (БС):

\begin{equation}

\label{krutov:lor}

 \vec{\tilde{j^0}} = \vec{j^0}+\frac{\vec{w}(\vec{j^0}\,\vec{w})}{1+w_{0}}-\vec{w}\,j^0_{0}, \quad

 \tilde{j^0}_{0} = j^0_{0}\,w_{0}-(\vec{j^0}\,\vec{w}) = j^0_{\mu}\,w^{\mu}\;,

\end{equation}

где $w_{\mu}={K'_{\mu}}/{\sqrt{K'^{2}}}$ "--- 4-скорость, соответствующая указанному преобразованию Лоренца;

$\tilde{j^0}_{\mu}$ "--- 4-вектор оператора электромагнитного тока в~БС.





Связь между матричными элементами токов в лабораторной системе \eqref{krutov:nodia} и системе Брейта  имеет следующий вид:

\begin{multline}

\label{krutov:diagonal111}

\langle\vec{W^{i}},\sqrt{s}|j_{\mu}^{0}(0)|\vec{W'^{i}},\sqrt{s'}, 1, 0, 1, m'\rangle =

\sum_{\tilde{m},\tilde{m}'}D^{0}_{0,\tilde{m}}(W^{i},w)D^{*{1}}_{m',\tilde{m}'}(W'^{i},w) \times\\

\times\langle\vec{\tilde{W^{i}}},\sqrt{s}|\tilde{j}_{\mu}^{0}(0)|\vec{\tilde{W'^{i}}},\sqrt{s'}, 1, 0, 1,\tilde{m}'\rangle,

\end{multline}

\begin{gather}

\tilde{W^{i}}=(\tilde{W^{i}}_{0},\vec{q}), \quad \tilde{W'^{i}}=(\tilde{W'^{i}}_{0},-\vec{q}),

\\

\tilde{W^{i}}_{0}=\sqrt{s+(\vec{q})^{2}}, \quad \tilde{W'^{i}}_{0}=\sqrt{s'+(\vec{q})^{2}} \quad

K'_{\mu}=(\sqrt{K'^{2}}, 0, 0, 0),

\label{krutov:breit}

\end{gather}

где $\vec{q}$ "--- трехмерный вектор с модулем

$$q=\sqrt{\lambda(M^2_{1},M^2_{2},Q^2)/[8(M_{1}^{2}+M_{2}^{2})+4Q^2]},$$

 $Q^2$ "--- квадрат переданного импульса.



Раскладывая нулевую компоненту матричного элемента оператора тока

в~БС \eqref{krutov:diagonal111} по сферическим углам вектора $\vec{q}$ и~применяя теорему Вигнера"--~Эккарта \cite{Zar93}, получим

\begin{multline}\label{krutov:A11}

\langle\vec{W^{i}},\sqrt{s}|{j_0}(0)|\vec{W'^{i}},\sqrt{s'}, 1, 0, 1, m'\rangle=

\\

=

\sum_{ \tilde{m}', \tilde{m}, l', k'}D^{0}_{0

\tilde{m}}(W^{i},w) D^{1}_{m'\tilde{m}'}(W'^{i},w)

\times 

\\

\times

\langle 1 \tilde{m}' l'k' |0 \tilde{m}\rangle \cdot Y_{l' k'}(\vec{q})\cdot G^{0, l'}_{01}(s,Q^2,s'),

\end{multline}

где $G^{0, l'}_{01}(s,Q^2,s')$ представляет собой набор скаляров или свободных электромагнитных формфакторов.



Перейдем теперь к рассмотрению трехмерной части матричного элемента оператора тока

\begin{equation}\label{krutov:breit1}

\langle\vec{{W^{i}}},\sqrt{s}|{{j^0}}_{t}(0)|\vec{{W'^{i}}},\sqrt{s'}, 1, 0, 1, {m}'\rangle,

\quad t = 1, 2, 3.

\end{equation}

Трехмерную часть оператора мы можем описать в терминах тензорного

оператора первого ранга. Для этого достаточно перейти к~каноническому базису, т.е. от декартовых компонент тока к~базису

сферических гармоник \cite{Edmons58}.



%\begin{equation}\label{krutov:B}

%{{j^0}}_{t}(0)=a_{t M}{{j^0}}^{1}_M(0), \quad M=-1,0,1\;,

%\end{equation}

%\begin{equation}\label{krutov:C}

%a_{t M}=\sqrt{\frac{2\pi}{3}}\left(

%\begin{tabular}{ccc}

%-1 & 0& 1\\

%i & 0& i\\

%0 & ${\sqrt{2}}$& 0\\

%\end{tabular}\right)\;,

%\end{equation}

%${{j^0}}^{1}_M(0)$ - компоненты тензорного оператора первого

%ранга.



Раскладывая матричный элемент \eqref{krutov:breit1} по сферическим углам

вектора $\vec{q}$ в~БС и применяя теорему Вигнера"--~Эккарта,

получим

\begin{multline}\label{krutov:Atok}

\langle\vec{W^{i}},\sqrt{s}|{j^0}^{1}_M(0)|\vec{W'^{i}},\sqrt{s'}, 1, 0, 1, m'\rangle=

\\

=

\sum_{\tilde{m}', \tilde{m}, M,k,l,j,n}D^{0}_{0 \tilde{m}}(W^{i},w)

D^{1}_{m'\tilde{m}'}(W'^{i},w)

\times \\

\times \langle 1\tilde{m}' j n |0

\tilde{m}\rangle \cdot \langle 1 M l k |j

n\rangle \cdot Y_{l k}(\vec{q})\cdot

G^{1,l,j}_{01}(s,Q^2,s').

\end{multline}



С другой стороны, матричный элемент тока \eqref{krutov:nodia} можно

записать в базисе индивидуальных переменных частиц системы:



\begin{multline}\label{krutov:Clebsh}

\langle\vec{W^i},\sqrt{s}|j^{0}_\mu(0)|\vec{W'^{i}},\sqrt{s'}, 1, 0, 1, m'\rangle =

\\

=\int\frac{d^3{\vec{w^i_1}}}{2w^i_{10}}\int\frac{d^3{\vec{w^i_2}}}{2w^i_{20}}\int\frac{d^3{\vec{w^{i}_1}'}}{2w^{i'}_{10}}\int

\frac{d^3{\vec{w^{i}_2}'}}{2w^{i'}_{20}}

\times \\

\\

\times

\langle\vec{W^i},\sqrt{s}|\vec{w^{i}_1},m_{1};\vec{w^{i}_2},m_{2}\rangle\cdot

\langle\vec{w^{i}_1},m_{1};\vec{w^{i}_2},m_{2}|j^{0}_\mu(0)|\vec{w^{i}_1}',m'_{1};\vec{w^{i}_2}',m'_{2}\rangle

\times

\\

\times 

\langle\vec{w^{i}_1}',m'_{1};\vec{w^{i}_2}',m'_{2}|\vec{W'^{i}},\sqrt{s'}, 1, 0, 1, m'\rangle\;,

\end{multline}

где

\begin{multline}\label{krutov:n1}

\langle\vec{w^i_1},m_{1};\vec{w^i_2},m_{2}|j^{0}_\mu(0)|\vec{w^i_1}',m'_{1};\vec{w^i_2}',m'_{2}\rangle

=\\

= \langle\vec{w^i_1},m_{1}|j^{0}_{1\mu }(0)|\vec{w^i_1}',m'_{1}\rangle \cdot 

\delta(\vec{w^i_2}-\vec{w^i_2}')\delta_{m_{2} m_{2'}} + \\

+

\langle\vec{w^i_2},m_{2}|j^{0}_{2\mu }(0)|\vec{w^i_2}',m'_{2}\rangle

\delta(\vec{w^i_1}-\vec{w^i_1}') \delta_{m_{1} m_{1'}}.

\end{multline}



Приравнивая покомпонентно выражение \eqref{krutov:Clebsh} с \eqref{krutov:A11} и \eqref{krutov:Atok}, используя разложение по каноническому базису и выполняя

интегрирование в (БС) в системе отсчета \mbox{$\vec{q}=(0, 0, q)$,} получим

аналитические выражения для свободных двухчастичных формфакторов.

В силу громоздкости данных выражений выпишем только один

формфактор, который в дальнейшем будет использоваться:

\begin{multline}\label{krutov:sv}

G_{01}^{111}(s,Q^2,s') =

\frac{\sqrt{2}\cdot\Theta(s,Q^2,s')(s+s'+Q^2)^{2}}{2\sqrt{s-4M^2}\sqrt{s'-4M^2}\sqrt{4M^2+Q^2}[\lambda(s,-Q^2,s')]^{1/2}}

\times

\\

\times

{\cos[({\omega_1}+{\omega_2})/2]}

\Bigl[\frac{s'(s'-s+3Q^2)}{[\lambda(s,-Q^2,s')]^{1/2}}\bigl(G_{M}^{u}(Q^2)+G_{M}^{\bar{d}}(Q^2)\bigr)\Bigr]+

\\

+{\sin[({\omega_1}+{\omega_2})/2]}

\Bigl[\frac{s'-s-Q^2}{s+s'+Q^2} \frac{\xi(s,s',Q^2)}{\sqrt{s'}}\bigl(G_{M}^{u}(Q^2)+G_{M}^{\bar{d}}(Q^2)\bigr)\Bigr]- 

\\

-{\sin[({\omega_1}+{\omega_2})/2]}

\Bigl[\xi(s,s',Q^2)\frac{4M}{s+s'+Q^2}\bigl(G_{E}^{u}(Q^2)+G_{E}^{\bar{d}}(Q^2)\bigr)

\Bigr],

\end{multline}

где

$$\Theta(s,Q^2,s')=\vartheta(s'-s_{1})-\vartheta(s'-s_{2}),~~ \xi(s,s',Q^2)=\sqrt{ss'Q^2-M^2\lambda(s,-Q^2,s')},$$

$$s_{1,2}=2M^{2}+\frac{1}{2M^2}(2M^2+Q^2)(s-2M^2)\mp\frac{1}{2M^2}\sqrt{Q^{2}(Q^{2}+4M^{2})s(s-4M^2)},$$

%$$\xi(s,s',Q^2)=\sqrt{ss'Q^2-M^2\lambda(s,-Q^2,s')},$$

$$\omega_1=\arctan\frac{\xi(s,s',Q^2)}{M[(\sqrt{s}+\sqrt{s'})^2+Q^2]+\sqrt{ss'}(\sqrt{s}+\sqrt{s'})},$$

$$\omega_2=\arctan\frac{\xi(s,s',Q^2)(2M+\sqrt{s}+\sqrt{s'})}{M(s+s'+Q^2)(2M+\sqrt{s}+\sqrt{s'})+\sqrt{ss'}(4M^2+Q^2)},$$

 $\vartheta$ "--- ступенчатая функция.



Отметим, что аналитическое выражение \eqref{krutov:sv} полностью совпадает с~перечисленными выше основными формами ПИКМ. 



\smallskip



\Section{Построение матричного элемента тока составной системы}

Матричный элемент электромагнитного тока перехода $\rho \to \pi\gamma^{*}$ для основных

форм ПИКМ можно записать следующим образом \cite{CaG95}:

\begin{equation}\label{krutov:tok}

\langle{W^i_{\pi}}|j_{\mu c}(0)|{W^i_{\rho}}, 1, m_{\rho}\rangle=n_{c}^{i}F_{\pi\rho}(Q^2)

\varepsilon_{\mu \nu \sigma \delta}

\eta^{\nu}(m_{\rho})W_{\pi}^{i\sigma}W_{\rho}^{i\delta},

\end{equation}

где $W^i_{\pi}$ и $W^i_{\rho}$ "--- 4-векторы $\pi$- и  $\rho$-мезонов для основных форм ПИКМ; 

$\eta^{\nu}(m_{\rho})$ "--- 4-вектор  поляризации;

$\varepsilon_{\mu \nu \sigma \delta}$ "--- антисимметричный тензор четвертого ранга; 

 $F_{\pi\rho}(Q^2)$ "--- формфактор, соответствующий  данному переходу; 

 $n_{c}^{i}$ "--- нормировочный множитель  разный для разных форм ПИКМ

$(n^{1}_{c}=1$, $n^{2}_{c}=1/M_{\pi}M_{\rho}$, $n^{3}_{c}=1)$; 

$M_{\pi}$ и $M_{\rho}$ "--- массы $\pi$- и $\rho$-мезонов.



Для дальнейшей работы с матричным элементом \eqref{krutov:tok} перейдем в

брейтовскую систему отсчета \eqref{krutov:breit}. 

В~выбранной системе отсчета вектор поляризации имеет следующий вид:

\begin{equation}\label{krutov:pol}

\eta^{\nu}(m_{\rho})=-\frac{1}{\sqrt{2}}(0, 0, 1, i).

\end{equation}





В системе Брейта выражение \eqref{krutov:tok} имеет вид

%

%Подставляя \eqref{krutov:A11}, \eqref{krutov:Atok} в~\eqref{krutov:tok} в~системе

%$\vec{q}=(0, 0, q),$ получим

\begin{equation}\label{krutov:tok1}

\langle{W^i{\pi}}|j_1(0)|{W^i_{\rho}}, 1, m_{\rho}\rangle=-\frac{q

\left(\sqrt{M_{\pi}^{2}+q^2}+\sqrt{M_{\rho}^{2}+q^2}\right)}{\sqrt{2}}F_{\pi\rho}(Q^2).

\end{equation}

Следует отметить, что компоненты матричного элемента тока \eqref{krutov:tok} $j_{0}$, $j_{2}$,  $j_{3}$ и соответствующие формфакторы оказываются равными нулю в ходе математических преобразований~\eqref{krutov:breit}. 

Таким образом, для радиационного перехода $\rho \to \pi\gamma^{*}$ существует единственный формфактор,

выраженный через первую трехмерную компоненту матричного элемента электромагнитного тока:

\begin{equation}\label{krutov:tok2}

\langle{W^i_{\pi}}|j_1(0)|{W^i_{\rho}}, 1, m_{\rho}\rangle=-\frac{1

}{\sqrt{3}}G_{01}^{111}(Q^2),

\end{equation}

где 

$G_{01}^{111}(Q^2)$ "--- формфактор составной системы, полученный при недиагональной параметризации.



Приравняв выражения \eqref{krutov:tok1} и \eqref{krutov:tok2}, получим

\begin{equation}

\label{krutov:transition}

F_{\pi\rho}(Q^2)=\sqrt{\frac{2}{3}} 

\frac{G_{01}^{111}(Q^2)}{q\bigl(\sqrt{M_{\pi}^{2}+q^2}+\sqrt{M_{\rho}^{2}+q^2}\bigr)},

\end{equation}



Используя модифицированное импульсное приближение \cite{KrT09}, можно показать, что

\begin{equation}\label{krutov:Fc1}

G^{1, l, 1}_{01}(Q^2)=\iint

G^{1, l, 1}_{01}(s,Q^2,s')\varphi(s)\varphi^{J'}_{S'}(s') \,d\sqrt{s} \, d\sqrt{s'},

\end{equation}

где

\begin{equation}

%\label{krutov:1}

\varphi(s)= \sqrt[4]{s}k\psi({k}),  \quad \varphi^{J'}_{S'}(s')= \sqrt[4]{s}k'\psi({k'});

\end{equation}

$\psi({k})$ и $\psi({k'})$ "--- волновые функции, удовлетворяющие условию

нормировки

\begin{equation}\label{krutov:q1}

\int_{-\infty}^{+\infty}\psi^2(k)k^{2}dk=1.

\end{equation}

Поставляя \eqref{krutov:Fc1} в \eqref{krutov:transition}, получим окончательное выражение для расчета переходного

формфактора $F_{\pi\rho}(Q^2)$:

\begin{multline}\label{krutov:transition1}

F_{\pi\rho}(Q^2)=\sqrt{\frac{2}{3}} 

\frac{1}{q\bigl(\sqrt{M_{\pi}^{2}+q^2}+\sqrt{M_{\rho}^{2}+q^2}\bigr)} 

\times 

\\

\times

\iint  G_{01}^{111}(s,Q^2,s')\varphi(s)\varphi^{J'}_{S'}(s') \, d\sqrt{s} \, d\sqrt{s'}.

\end{multline}



\smallskip



\Section{Расчет переходного формфактора $\boldsymbol{ F_{\pi \rho}(Q^2)}$}

Для расчета формфактора \eqref{krutov:transition1}

используем волновую функцию основного состояния гармонического осциллятора и волновую функцию степенного типа \cite{AKr11}:

\begin{gather}\label{krutov:volna1}

\psi(k)=\frac{2}{\pi^{1/4}a^{3/2}}\exp\Bigl(-\frac{k^2}{2a^2}\Bigr),

\\

\label{krutov:stepen}

\psi(k)=\sqrt{\frac{2^9}{7\pi b^{3}}}\frac{1}{(1 + k^2/b^2)^3},

\end{gather}

где $a$ и $b$ "--- параметры волновых функций, которые фиксируются из требования описания среднеквадратичных радиусов $\pi$- и $\rho$-мезонов \cite{KrTr03}.



Для расчета переходного формфактора выберем саксовские формфакторы кварков в виде

\begin{equation}\label{krutov:mmm}

G_{E}(Q^2)=e_{q}f_{q}(Q^2),\quad G_{M}(Q^2)=(e_{q}+\kappa_{q})f_{q}(Q^2),

\end{equation}

где $e_{q}$ "--- заряд кварка, $\kappa_{q}$ "--- аномальный магнитный момент кварка в~естественных единицах \cite{G95}. 

Для расчета функции $f_{q}(Q^2)$ будем использовать выражение, введенное в \cite{CaGNS94}:

\begin{equation}\label{krutov:f2}

f_{q}(Q^2)=\frac{1}{1+\langle r^2_{q}\rangle Q^2/6},

\end{equation}

где

$\langle r^2_{q}\rangle=0.3/M^2$ "--- среднеквадратичный радиус кварка~\cite{CaGNS94}.



Результаты расчета переходного формфактора \eqref{krutov:transition1}

представлены на рисунке.

\begin{figure}[h]

\footnotesize \centering

\centering

\includegraphics[width=75mm]{politex}

\caption*{Результаты расчета переходного формфактора (\ref{krutov:transition1})

 для разных типов волновых функций:   сплошная линия "--- расчет с использованием волновой функции (\ref{krutov:volna1});  штриховая линия "--- расчет с использованием волновой функции (\ref{krutov:stepen}) (online в цвете)}



\smallskip



[The results of the calculation of the transition form factor (\ref{krutov:transition1}) for different types of wave functions. Solid line is the results of the calculation using the wave function \eqref{krutov:volna1}; 

dashed~line~is~the results of a calculation using the wave function \eqref{krutov:stepen}  (color online)]  



 

\end{figure}





\Section[n]{Заключение}

В рамках подхода ПИКМ проведена процедура построения матричного

элемента тока перехода  недиагонального по полному угловому

моменту. Изложение данной методики проводилось на примере

радиационного перехода $\rho \to \pi\gamma^{*}$. Для решения

задачи  в модифицированном релятивистском импульсном приближении

проведено построение оператора электромагнитного тока для

составной двухчастичной системы с учетом условий

лоренц-ковариантности и сохранения. Получено аналитическое

выражение для формфактора указанного перехода, которое совпадает

с основными формами ПИКМ. Произведен численный расчет переходного

формфактора с  волновыми функциями двух видов.









\thanks{{\bf Благодарности.} Авторы выражают благодарность В.~Е.~Троицкому  за плодотворные обсуждения.

}



\thanks{{\bf ORCIDs}



{\bf Александр Федорович Крутов:} \url{http://orcid.org/0000-0002-8484-0028}



{\bf Роман Геннадьевич Полежаев:} \url{http://orcid.org/0000-0002-8781-4116}



}







\RusRef





\begin{thebibliography}{20}



\RBibitem{krutov:add1-ru}

\by  Крутов~А.~Ф., Полежаев~Р.~Г.

\paper Описание радиационных распадов $V \to  P\gamma^{*}$ в различных формах Пуанкаре-инвариантной квантовой механики

\pages 210--211

\inbook Четвертая международная конференция «Математическая физика и ее приложения»

\bookinfo материалы конф. 

\eds чл.-корр.~РАН И.~В.~Волович; д.ф.-м.н., проф. В.~П.~Радченко

\publaddr Самара

\publ СамГТУ

\yr 2014







\Bibitem{Na6009}

\by Arnaldi~R. (et.al. NA60 Collaboration)

\paper Study of the electromagnetic transition form-factors in $\eta\to \mu^{+}\mu^{-}\gamma$ and $\omega\to \mu^{+}\mu^{-}\pi^{0}$ decays with NA60

\jour Phys. Lett.~B

\yr 2009

\vol 677

\pages 260--266

\issue 5

\crossref{10.1016/j.physletb.2009.05.029}

\arxiv \href{http://arxiv.org/abs/0902.2547}{0902.2547} [hep-ph]



\Bibitem{Na6011}

\by Usai~G. (et.al. NA60 Collaboration)

\paper Low mass dimuon production in proton-nucleus collisions at 400~GeV/c

\jour Nucl. Phys.~A

\yr 2011

\vol 855

\issue 1

\pages 918--196

\crossref{10.1016/j.nuclphysa.2011.02.037}



\Bibitem{Na60Ur11}

\by Uras~A. (et.al. NA60 Collaboration)

\paper Measurement of the $\eta$ and $\omega$ Dalitz decays transition form factors in p-A collisions at 400~GeV/c with the NA60 apparatus

\jour J. Phys.: Conf. Ser. 

\yr 2011

\vol 270

\issue 1

\papernumber 012038

\crossref{10.1088/1742-6596/270/1/012038}

\arxiv \href{http://arxiv.org/abs/1108.0968}{1108.0968} [hep-ex]



\Bibitem{Kloe11}

\by Archilli~F. (et.al. KLOE-2 Collaboration)

\crossref{10.1016/j.physletb.2011.11.033}

\paper  Search for a vector gauge boson in $\phi$ meson decays with the KLOE detector

\jour Phys. Lett.~B

\vol 706

\issue 4--5

\arxiv \href{http://arxiv.org/abs/1110.0411}{1110.0411} [hep-ph]

\yr 2012

\pages  251--255



\Bibitem{Amsler08}

\by Amsler~C. (et.al. Particle Data Group)

\paper   Review of particle physics

\jour Phys. Lett.~B

\yr 2008

\vol 667

\issue 1--5

\pages 1--6

\crossref{10.1016/j.physletb.2008.07.018}







\RBibitem{KrT09}

\by Крутов~А.~Ф., Троицкий~В.~Е.

\paper   Мгновенная форма Пуанкаре-инвариантной квантовой механики и описание структуры составных систем

\jour Физика элементарных частиц и атомного ядра

\yr 2009

\vol 40

\issue 2

\pages 269--319

\elib{http://elibrary.ru/item.asp?id=13053810}

%Instant form of Poincar?-invariant quantum mechanics and description of the structure of composite systems

%Krutov A.F., Troitsky V.E.

%Physics of Particles and Nuclei. 2009. Т. 40. № 2. С. 136-161.

% 10.1134/S1063779609020026

%\elib{http://elibrary.ru/item.asp?id=13605887}





\Bibitem{MaT02}

\by Maris~P., Tandy~P.C.

\paper    Electromagnetic transition form factors of light mesons 

\jour Phys.~Rev.~C

\yr 2002

\vol 65

\issue 4

\papernumber 045211

\crossref{10.1103/physrevc.65.045211}

\arxiv \href{http://arxiv.org/abs/nucl-th/0201017}{nucl-th/0201017}



\Bibitem{JBol07}

\by Yu~J., Xiao~B.-W.,  Ma~B.-Q.

\paper  Space-like and time-like pion–rho transition form factors in the light-cone formalism 

\jour J.~Phys. G: Nucl. Part. Phys. 

\yr 2007

\vol 34

\issue 7

\pages 1845--1860

\arxiv \href{http://arxiv.org/abs/0706.2018}{0706.2018} [hep-ph]

\crossref{10.1088/0954-3899/34/7/021}



\Bibitem{Displ09}

\by Desplanques~B.

\paper   RQM description of the charge form factor of the pion and its asymptotic behavior

\jour Eur. Phys.~J.~A

\yr 2009

\vol 42

\issue 2

\pages 219--236

\arxiv \href{http://arxiv.org/abs/0906.1889}{0906.1889} [nucl-th]

\crossref{10.1140/epja/i2009-10864-8}





\Bibitem{Iva12}

\by Ivashyn~S.A.

\paper    Vector to pseudoscalar meson radiative transition in chiral theory with resonances

\inbook Problems of Atomic Science and Technology. \rm No.~1

\serial Nuclear Physics Investigations 

\yr 2012

\vol 57

\pages 179--182

\arxiv \href{http://arxiv.org/abs/1111.1291}{1111.1291} [hep-ph]



\Bibitem{Klink14}

\by  Bierrat~E.~P., Schweiger~W.

\paper  Electromagnetic $\rho$-meson form factors in point-form relativistic quantum mechanics

\jour  Phys.~Rev.~C

\yr 2014

\vol 89

\issue 5

\papernumber 055205

\crossref{10.1103/physrevc.89.055205}

\arxiv \href{http://arxiv.org/abs/1404.2440}{1404.2440} [hep-ph]







\Bibitem{KrTrTs09}

\by Krutov~A.~F., Troitsky~V.~E., Tsirova~N.~A.

\paper  Nonperturbative relativistic approach to pion form factor: predictions for future JLab experiments

\jour Phys.~Rev.~C

\yr 2009

\vol 80

\issue 5

\papernumber 055210

\crossref{10.1103/physrevc.80.055210}

\elib{http://elibrary.ru/item.asp?id=15305002}

\arxiv \href{http://arxiv.org/abs/0910.3604}{0910.3604} [nucl-th]







\RBibitem{KrPol13-2}

\by Крутов~А.~Ф., Полежаев~Р.~Г.

\paper  Описание электромагнитной структуры пиона в~различных формах Пуанкаре-инвариантной квантовой механики

\jour Ядерная физика и инжиниринг

\yr 2013

\vol 4

\issue 9--10

\pages 848--852

\elib{http://elibrary.ru/item.asp?id=20991932}

\crossref{10.1134/S2079562913090200}



\RBibitem{ShirCh63}

\by Чешков~А.~А.,  Широков~Ю.~М.

\paper   Инвариантная параметризация локальных операторов

\jour Ж.~экспер. теорет. физ.

\yr 1963

\vol 44

\pages 1982--1992

%MR0160501

%Cheshkov, A. A.; Shirokov, Yu. M.

%Invariant parametrization of local operators. ?. ?ksper. Teoret. Fiz. 44 1982--1992 (Russian. English summary); translated as

%Soviet Physics JETP 17 1963 1333–1339. 

% Invariant parametrization of local operators

%

%Journal Article published Nov 1963 in Nuclear Physics volume 49 on pages 108 to 120

%

%Authors: A.A. Cheshkov, Yu.M. Shirokov

%http://dx.doi.org/10.1016/0029-5582(63)90079-8 



\Bibitem{Zar93}

\by Zare~R.~N.

\book Angular Momentum: Understanding Spatial Aspects in Chemistry and Physics

\publaddr New York

\publ Wiley

\yr 1988

\totalpages xi+349



\Bibitem{Edmons58} 

\book  Angular Momentum in Quantum Mechanics

\yr 1957

\by Edmonds~A.~R.

\vol 4

\serial Investigations in Physics

\publaddr  Princeton, New Jersey

\publ Princeton University Press 

\totalpages viii+146 

\zmath{https://zbmath.org/?q=an:0079.42204}



\Bibitem{CaG95}

\by  Cardarelli~F., Grach~I.~L., Narodetskii~I.~M., Salm\'e~G., Simula~S.

\paper    Radiative $\pi\rho$ and $\pi\omega$ transition form factor in a light-front constituent quark model

\jour Phys. Lett.~B

\yr 1995

\vol 359

\issue 1--2

\pages 1--7

\crossref{10.1016/0370-2693(95)01058-x}

\arxiv \href{http://arxiv.org/abs/nucl-th/9509004}{nucl-th/9509004}



\RBibitem{AKr11}

\by  Андреев~В.~В., Крутов~А.~Ф.

\paper    Электромагнитные формфакторы мезонов

\jour Проблемы физики, математики и техники

\yr 2011

\issue 1(6)

\pages 7--19

\elib{http://elibrary.ru/item.asp?id=18925725}



\RBibitem{KrTr03}

\by Крутов~А.~Ф., Троицкий~В.~Е.

\paper   Релятивистские эффекты в~электромагнитной структуре $\rho$-мезона

\jour Вестн. СамГУ. Естественнонаучн. сер.

\yr 2003

\issueinfo Второй спец. выпуск

\pages 95--111



\RBibitem{G95}

\by  Gerasimov~S.~B.

\paper  Magnetic moments of baryons and strange content of the nucleon

\jour Phys. Lett.~B

\yr 1995

\vol 357

\pages 666--670

\issue 4

\crossref{10.1016/0370-2693(95)00934-d}





\Bibitem{CaGNS94}

\by   Cardarelli~F., Grach~I.~L., Narodetskii~I.~M., Pace~E., Salme~G., Simula~S.

\paper Hard Constituent Quarks and Electroweak Properties of Pseudoscalar Mesons

\jour Phys. Lett.~B

\yr 1994

\vol 332

\issue 1--2

\pages 1--7

\crossref{10.1016/0370-2693(94)90849-4}

\arxiv \href{http://arxiv.org/abs/nucl-th/9405014}{nucl-th/9405014}



\end{thebibliography}



\RuDate





\pagebreak



\makeentitle





\thanks{{\bf Acknowledgments.} The authors are grateful to Vadim E.~Troitsky  for helpful comments and insightful discussions.

}



\thanks{{\bf ORCIDs}



{\bf Alexander F. Krutov:} \url{http://orcid.org/0000-0002-8484-0028}



{\bf Roman G. Polezhaev:} \url{http://orcid.org/0000-0002-8781-4116}



}







\EngRef





\begin{thebibliography}{20}



\Bibitem{krutov:add1-eng}

\by  Krutov~A.~F., Polezhaev~R.~G.

\paper Description of radiative decays of $V \to P \gamma^{*}$ in different forms of Poincar\'e-invariant quantum mechanics

\pages 210--211

\inbook  The 4nd International Conference ``Mathematical Physics and its Applications''

\bookinfo Book of Abstracts and Conference Materials

\eds I.~V.~Volovich; V.~P.~Radchenko

\publaddr Samara

\publ Samara State Technical Univ.

\yr 2014

\lang In Russian





\Bibitem{Na6009:eng}

\by Arnaldi~R. (et.al. NA60 Collaboration)

\paper Study of the electromagnetic transition form-factors in $\eta\to \mu^{+}\mu^{-}\gamma$ and $\omega\to \mu^{+}\mu^{-}\pi^{0}$ decays with NA60

\jour Phys. Lett.~B

\yr 2009

\vol 677

\pages 260--266

\issue 5

\crossref{10.1016/j.physletb.2009.05.029}

\arxiv \href{http://arxiv.org/abs/0902.2547}{0902.2547} [hep-ph]



\Bibitem{Na6011:eng}

\by Usai~G. (et.al. NA60 Collaboration)

\paper Low mass dimuon production in proton-nucleus collisions at 400~GeV/c

\jour Nucl. Phys.~A

\yr 2011

\vol 855

\issue 1

\pages 918--196

\crossref{10.1016/j.nuclphysa.2011.02.037}



\Bibitem{Na60Ur11:eng}

\by Uras~A. (et.al. NA60 Collaboration)

\paper Measurement of the $\eta$ and $\omega$ Dalitz decays transition form factors in p-A collisions at 400~GeV/c with the NA60 apparatus

\jour J. Phys.: Conf. Ser. 

\yr 2011

\vol 270

\issue 1

\papernumber 012038

\crossref{10.1088/1742-6596/270/1/012038}

\arxiv \href{http://arxiv.org/abs/1108.0968}{1108.0968} [hep-ex]



\Bibitem{Kloe11:eng}

\by Archilli~F. (et.al. KLOE-2 Collaboration)

\crossref{10.1016/j.physletb.2011.11.033}

\paper  Search for a vector gauge boson in $\phi$ meson decays with the KLOE detector

\jour Phys. Lett.~B

\vol 706

\issue 4--5

\arxiv \href{http://arxiv.org/abs/1110.0411}{1110.0411} [hep-ph]

\yr 2012

\pages  251--255



\Bibitem{Amsler08:eng}

\by Amsler~C. (et.al. Particle Data Group)

\paper   Review of particle physics

\jour Phys. Lett.~B

\yr 2008

\vol 667

\issue 1--5

\pages 1--6

\crossref{10.1016/j.physletb.2008.07.018}







\Bibitem{KrT09:eng}

\paper Instant form of Poincar\'e-invariant quantum mechanics and description of the structure of composite systems

\by Krutov~A.~F., Troitsky~V.~E.

\jour Physics of Particles and Nuclei

\yr 2009

\vol 40

\issue 2

\pages 136--161

\crossref{10.1134/S1063779609020026}

\elib{http://elibrary.ru/item.asp?id=13605887}





\Bibitem{MaT02:eng}

\by Maris~P., Tandy~P.C.

\paper    Electromagnetic transition form factors of light mesons 

\jour Phys.~Rev.~C

\yr 2002

\vol 65

\issue 4

\papernumber 045211

\crossref{10.1103/physrevc.65.045211}

\arxiv \href{http://arxiv.org/abs/nucl-th/0201017}{nucl-th/0201017}



\Bibitem{JBol07:eng}

\by Yu~J., Xiao~B.-W.,  Ma~B.-Q.

\paper  Space-like and time-like pion–rho transition form factors in the light-cone formalism 

\jour J.~Phys. G: Nucl. Part. Phys. 

\yr 2007

\vol 34

\issue 7

\pages 1845--1860

\arxiv \href{http://arxiv.org/abs/0706.2018}{0706.2018} [hep-ph]

\crossref{10.1088/0954-3899/34/7/021}



\Bibitem{Displ09:eng}

\by Desplanques~B.

\paper   RQM description of the charge form factor of the pion and its asymptotic behavior

\jour Eur. Phys.~J.~A

\yr 2009

\vol 42

\issue 2

\pages 219--236

\arxiv \href{http://arxiv.org/abs/0906.1889}{0906.1889} [nucl-th]

\crossref{10.1140/epja/i2009-10864-8}





\Bibitem{Iva12:eng}

\by Ivashyn~S.A.

\paper    Vector to pseudoscalar meson radiative transition in chiral theory with resonances

\inbook Problems of Atomic Science and Technology. \rm No.~1

\serial Nuclear Physics Investigations 

\yr 2012

\vol 57

\pages 179--182

\arxiv \href{http://arxiv.org/abs/1111.1291}{1111.1291} [hep-ph]



\Bibitem{Klink14:eng}

\by  Bierrat~E.~P., Schweiger~W.

\paper  Electromagnetic $\rho$-meson form factors in point-form relativistic quantum mechanics

\jour  Phys.~Rev.~C

\yr 2014

\vol 89

\issue 5

\papernumber 055205

\crossref{10.1103/physrevc.89.055205}

\arxiv \href{http://arxiv.org/abs/1404.2440}{1404.2440} [hep-ph]







\Bibitem{KrTrTs09:eng}

\by Krutov~A.~F., Troitsky~V.~E., Tsirova~N.~A.

\paper  Nonperturbative relativistic approach to pion form factor: predictions for future JLab experiments

\jour Phys.~Rev.~C

\yr 2009

\vol 80

\issue 5

\papernumber 055210

\crossref{10.1103/physrevc.80.055210}

\elib{http://elibrary.ru/item.asp?id=15305002}

\arxiv \href{http://arxiv.org/abs/0910.3604}{0910.3604} [nucl-th]







\Bibitem{KrPol13-2:eng}

\lang In Russian

\by  Krutov~A.~F., Polezhaev~R.~G.

\paper  Description of the electromagnetic structure of the pion in the various forms of the Poincar\'e-invariant quantum mechanics

\jour Iadernaia fizika i inzhiniring \rm [Nuclear Physics and Engineering]

\yr 2013

\vol 4

\issue 9--10

\pages 848--852

\crossref{10.1134/S2079562913090200}



\Bibitem{ShirCh63:eng}

%MR0160501

\by Cheshkov~A.~A., Shirokov~Yu.~M.

\paper Invariant parametrization of local operators

\jour Soviet Physics JETP 

\vol 17 

\yr 1963 

\pages 1333–1339

\moreref

\by Cheshkov~A.~A., Shirokov~Yu.~M.

\paper Invariant parametrization of local operators

\yr 1963 

\jour Nuclear Physics 

\vol 49 

\pages 108--120

\crossref{10.1016/0029-5582(63)90079-8} 



\Bibitem{Zar93:eng}

\by Zare~R.~N.

\book Angular Momentum: Understanding Spatial Aspects in Chemistry and Physics

\publaddr New York

\publ Wiley

\yr 1988

\totalpages xi+349



\Bibitem{Edmons58:eng} 

\book  Angular Momentum in Quantum Mechanics

\yr 1957

\by Edmonds~A.~R.

\vol 4

\serial Investigations in Physics

\publaddr  Princeton, New Jersey

\publ Princeton University Press 

\totalpages viii+146 

\zmath{https://zbmath.org/?q=an:0079.42204}



\Bibitem{CaG95:eng}

\by  Cardarelli~F., Grach~I.~L., Narodetskii~I.~M., Salm\'e~G., Simula~S.

\paper    Radiative $\pi\rho$ and $\pi\omega$ transition form factor in a light-front constituent quark model

\jour Phys. Lett.~B

\yr 1995

\vol 359

\issue 1--2

\pages 1--7

\crossref{10.1016/0370-2693(95)01058-x}

\arxiv \href{http://arxiv.org/abs/nucl-th/9509004}{nucl-th/9509004}



\Bibitem{AKr11:eng}

\lang In Russian

\by  Andreev~V.~V., Krutov~A.~F.

\paper    Electromagnetic form factors of mesons

\jour Problemy fiziki, matematiki i tekhniki

\yr 2011

\issue 1(6)

\pages 7--19





\Bibitem{KrTr03:eng}

\lang In Russian

\by Krutov~A.~F., Troitsky~V.~E.

\paper   Relativistic effects in the electromagnetic structure of the $\rho$-meson

\jour Vestnik SamGU. Estestvenno-Nauchnaya Ser.

\yr 2003

\issueinfo Second Special Issue

\pages 95--111



\RBibitem{G95:eng}

\by  Gerasimov~S.~B.

\paper  Magnetic moments of baryons and strange content of the nucleon

\jour Phys. Lett.~B

\yr 1995

\vol 357

\pages 666--670

\issue 4

\crossref{10.1016/0370-2693(95)00934-d}





\Bibitem{CaGNS94:eng}

\by   Cardarelli~F., Grach~I.~L., Narodetskii~I.~M., Pace~E., Salme~G., Simula~S.

\paper Hard Constituent Quarks and Electroweak Properties of Pseudoscalar Mesons

\jour Phys. Lett.~B

\yr 1994

\vol 332

\issue 1--2

\pages 1--7

\crossref{10.1016/0370-2693(94)90849-4}

\arxiv \href{http://arxiv.org/abs/nucl-th/9405014}{nucl-th/9405014}



\end{thebibliography}



\EngDate







\renewcommand{\baselinestretch}{0.9}





\Russian

\newpage



%%

%\end{document}

